\documentclass[a4paper,11pt]{article}
\usepackage[T1]{fontenc}
\usepackage[utf8]{inputenc}
\usepackage{lmodern}
\usepackage[indonesian]{babel}
\usepackage[margin=2.5cm]{geometry}
\usepackage{amsmath,amssymb,amsthm,mathtools}
\usepackage{graphicx}
\usepackage{enumitem}
\usepackage{microtype}
\usepackage[hidelinks]{hyperref}
\graphicspath{{gambar/}}
\setlength{\parindent}{1.25em}
\setlength{\parskip}{0.3em}
\setlength{\emergencystretch}{1em}
\raggedbottom
\setlist[enumerate]{label=\arabic*.,itemsep=0.3em,topsep=0.4em}
\addto\captionsindonesian{%
  \renewcommand{\abstractname}{Abstrak}%
  \renewcommand{\refname}{Daftar Pustaka}%
  \renewcommand{\proofname}{Bukti}%
  \renewcommand{\figurename}{Gambar}}

% Satu urutan nomor 1--39, sama dengan artikel sumber.
% Pernyataan, notasi, dan ketidakkonsistenan sumber dipertahankan.
% Tidak disertakan koreksi, asumsi tambahan, atau contoh buatan pengulas.
\theoremstyle{plain}
\newtheorem{definisi}{Definisi}
\newtheorem{contoh}[definisi]{Contoh}
\newtheorem{catatan}[definisi]{Catatan}
\newtheorem{teorema}[definisi]{Teorema}
\newtheorem{proposisi}[definisi]{Proposisi}
\newtheorem{lema}[definisi]{Lema}
\newtheorem{akibat}[definisi]{Akibat}
\newcommand{\F}{\mathbb{F}}
\newcommand{\R}{\mathbb{R}}
\newcommand{\Z}{\mathbb{Z}}
\newcommand{\upp}[1]{\overline{App}(#1)}
\newcommand{\lowp}[1]{\underline{App}(#1)}
\newcommand{\upr}[1]{\overline{Apr}(#1)}
\newcommand{\lowr}[1]{\underline{Apr}(#1)}
\newcommand{\res}[1]{[\![#1]\!]_3}

\title{\textbf{Ruang Vektor Rough}\\[0.4em]
\large}
\author{\textit{Journal: A New Approach to Rough Vector Spaces}\\
Author: Shahad T. Almohammadi \& Cenap \"Ozel}
\date{}

\begin{document}
\maketitle
\section{Konsep Dasar}

\begin{definisi}\label{def:1}
      Pasangan $(U,R)$ yang terdiri atas himpunan tak kosong $U$ dan relasi
      ekuivalensi $R$ disebut \textbf{ruang aproksimasi}.
\end{definisi}

\begin{definisi}\label{def:2}
      Misalkan $X$ sebarang himpunan bagian dari ruang aproksimasi $(U,R)$.
      Maka $X$ dinyatakan oleh pasangan $(\lowp{X},\upp{X})$, yang
      berturut-turut merupakan aproksimasi bawah dan aproksimasi atas, dengan
      \begin{align*}
            \lowp{X} & =\bigcup_{x\in U}\{[x]_R:[x]_R\subseteq X\},\quad\text{dan} \\
            \upp{X}  & =\bigcup_{x\in U}\{[x]_R:[x]_R\cap X\ne\varnothing\},
      \end{align*}
      di mana $[x]_R$ menyatakan kelas ekuivalensi yang memuat $x$ dalam $R$.
      \textbf{Batas} himpunan $X$ dalam $A$ didefinisikan sebagai
      \[
            Bnd(X)=\upr{X}-\lowr{X}.
      \]
\end{definisi}

Misalkan $U$ suatu ruang aproksimasi dan $X,Y\subseteq U$. Maka berlaku:
\begin{enumerate}
      \item $\lowp{X}\subseteq X\subseteq\upp{X}$,
      \item $\upp{U}=\lowp{U}=U$,
      \item $\upp{\phi}=\lowp{\phi}=\phi$,
      \item $\upp{X\cup Y}=\upp{X}\cup\upp{Y}$,
      \item $\lowp{X\cap Y}=\lowp{X}\cap\lowp{Y}$,
      \item $\upp{X\cap Y}\subset\upp{X}\cap\upp{Y}$,
      \item $\lowp{X}\cup\lowp{Y}\subseteq\lowp{X\cup Y}$,
      \item $\upp{X}\subset\upp{Y}$ dan $\lowp{X}\subset\lowp{Y}$
            apabila $X\subset Y$.
\end{enumerate}

\begin{definisi}\label{def:4}
      Himpunan bagian $G$ dari ruang aproksimasi $(U,R)$ dengan operasi biner
      $(*)$ yang didefinisikan pada $U$ disebut \textbf{grup rough} apabila
      empat aksioma berikut berlaku:
      \begin{enumerate}
            \item $x*y\in\upp{G}$, $\forall x,y\in G$;
            \item sifat asosiatif berlaku pada $\upp{G}$;
            \item $\exists e\in\upp{G}$ sehingga $\forall x\in G$,
                  $x*e=e*x=x$, dengan $e$ disebut identitas rough; dan
            \item $\forall x\in G$, $\exists y\in G$ sehingga $x*y=y*x=e$,
                  dengan $y$ disebut invers rough dan dinotasikan dengan $x^{-1}$.
      \end{enumerate}
\end{definisi}

\begin{contoh}\label{ex:5}
      Misalkan $X=\Z_3\times\Z_3=\{\langle a,b\rangle:a,b\in\Z_3\}$
      dengan operasi penjumlahan modulo $3$, dan misalkan
      \[
            V=\{\langle1,1\rangle,\langle1,2\rangle,\langle2,2\rangle,
            \langle2,1\rangle,\langle1,0\rangle,\langle2,0\rangle\}\subseteq X.
      \]
      Klasifikasi $X$ adalah $X/R=\{C_1,C_2\}$, dengan $C_1$ merupakan
      himpunan semua pasangan terurut yang sekurang-kurangnya satu
      komponennya nol, dan $C_2=X\setminus C_1$. Maka $\upp{V}=X$.

      Mari kita periksa bahwa $V$ merupakan grup rough. Berdasarkan
      Definisi~\ref{def:4}, syarat-syarat berikut terpenuhi:
      \begin{enumerate}
            \item $x*y\in\upp{V}$, $\forall x,y\in V$, karena $\upp{V}=X$;
            \item sifat asosiatif sudah berlaku pada $\upp{V}$;
            \item $\langle0,0\rangle\in\upp{V}$ dianggap sebagai identitas rough;
            \item $\langle1,1\rangle^{-1}=\langle2,2\rangle\in V$,
                  $\langle1,2\rangle^{-1}=\langle2,1\rangle\in V$, dan
                  $\langle1,0\rangle^{-1}=\langle2,0\rangle\in V$.
      \end{enumerate}
      Dengan demikian, $V$ merupakan grup rough.
\end{contoh}

\begin{contoh}[Pengaruh relasi ekuivalensi]\label{ex:relasi-grup}
      Ambil $U=\Z_6$ dengan penjumlahan modulo $6$ dan $G=\{1,5\}$.
      Himpunan $G$ bukan grup biasa karena $1+1=2\notin G$.
      Pertimbangkan dua relasi ekuivalensi berikut.
      \begin{enumerate}
            \item Jika $U/R_1=\{\{0\},\{1\},\{2\},\{3\},\{4\},\{5\}\}$,
                  maka $\overline{App}_{R_1}(G)=\{1,5\}$.
                  Karena $1+1=2\notin\overline{App}_{R_1}(G)$,
                  $G$ bukan grup rough terhadap $(U,R_1)$.

            \item Jika $U/R_2=\{\{0,1,2\},\{4,5\},\{3\}\}$, maka
                  \[
                        \overline{App}_{R_2}(G)=\{0,1,2,4,5\}.
                  \]
                  Semua kemungkinan jumlah dua anggota $G$ adalah
                  \[
                        1+1=2,\qquad 1+5=5+1=0,\qquad 5+5=4,
                  \]
                  yang seluruhnya berada di $\overline{App}_{R_2}(G)$.
                  Asosiativitas diwarisi dari penjumlahan pada $\Z_6$.
                  Elemen $0\in\overline{App}_{R_2}(G)$ merupakan identitas,
                  sedangkan invers aditif $1$ dan $5$ masing-masing adalah
                  $5$ dan $1$, yang keduanya berada di $G$.
                  Jadi, $G$ merupakan grup rough terhadap $(U,R_2)$.
      \end{enumerate}
      Dengan demikian, perubahan relasi ekuivalensi dapat mengubah
      status $G$ sebagai grup rough. Pada kasus kedua, berlaku
      \[
            G\subsetneq\overline{App}_{R_2}(G)\subsetneq U,
            \qquad \underline{App}_{R_2}(G)=\varnothing.
      \]
\end{contoh}


\begin{definisi}\label{def:8}
      Jika $G$ merupakan grup rough dan $H\subseteq G$, maka $H$ disebut
      \textbf{subgrup rough} dari $G$ apabila $H$ sendiri merupakan grup
      rough dengan operasi yang diinduksi dari $G$.
\end{definisi}

\begin{catatan}\label{rem:9}
      Satu-satunya subgrup rough trivial dari suatu grup rough $G$ adalah
      $G$ sendiri, dan himpunan bagian $\{e\}$ menjadi subgrup rough trivial
      hanya jika $e\in G$.
\end{catatan}

\begin{definisi}\label{def:10}
      Himpunan bagian $H$ dari grup rough $G$ merupakan subgrup rough dari
      $G$ jika dan hanya jika
      \begin{enumerate}
            \item $xy\in\upp{H}$, $\forall x,y\in H$;
            \item $x^{-1}\in H$, $\forall x\in H$.
      \end{enumerate}
\end{definisi}

\section{Ruang Vektor Rough}

Pada bagian ini, akan diberikan definisi ruang vektor rough beserta
sifat-sifat dan contoh-contoh yang berkaitan dengannya.

\begin{definisi}\label{def:15}
      Misalkan $(U,R)$ suatu ruang aproksimasi. Himpunan bagian $V\subseteq U$
      disebut \textbf{ruang vektor rough} atas lapangan $\F$, dinotasikan dengan
      $\boldsymbol{RV}_{\F}$, jika $V$ merupakan grup rough terhadap operasi
      biner $(+)$ dan sifat komutatif $(+)$ berlaku pada $\upp{V}$. Selain
      itu, untuk setiap $\lambda\in\F$ dan $v\in V$, terdapat elemen
      $\lambda\cdot v\in\upp{V}$ sehingga aksioma-aksioma berikut dipenuhi
      untuk setiap $\alpha,\beta\in\F$ dan setiap $u,v\in V$:
      \begin{enumerate}
            \item $\alpha\cdot(u+v)=\alpha\cdot u+\alpha\cdot v$;
            \item $(\alpha+\beta)v=\alpha\cdot v+\beta\cdot v$;
            \item $\alpha\cdot(\beta\cdot v)=(\alpha\cdot\beta)\cdot v$; dan
            \item $1\cdot v=v$, dengan $1\in\F$.
      \end{enumerate}
\end{definisi}

\begin{definisi}\label{def:16}
      Misalkan $A=(U,R)$ suatu ruang aproksimasi dan $V\subseteq U$ suatu
      ruang vektor rough atas lapangan $\F$. Himpunan bagian tak kosong
      $X\subset V$ disebut \textbf{subruang rough} dari $V$ atas lapangan
      $\F$ yang sama apabila $X$ sendiri merupakan ruang vektor rough
      terhadap operasi yang sama dengan operasi yang didefinisikan pada $V$.
\end{definisi}

\begin{teorema}\label{def:17}
      Himpunan bagian $W$ dari ruang vektor rough $V$ atas lapangan $\F$
      merupakan subruang rough jika dan hanya jika kedua syarat berikut
      terpenuhi:
      \begin{enumerate}
            \item $W$ merupakan subgrup rough;
            \item $\lambda x\in\upp{W}$, $\forall x\in W$ dan $\forall\lambda\in\F$.
      \end{enumerate}
\end{teorema}

\begin{contoh}\label{ex:18}
      Ambil lapangan $\Z_7$ dan ruang vektor $X=\Z_7^3$, dengan operasi
      penjumlahan dan perkalian skalar pada setiap koordinat modulo $7$.
      Tetapkan
      \[
            H=\{(a,b,0):a,b\in\Z_7\},\qquad
            V=\{(a,b,0):a\in\Z_7\setminus\{0\},\ b\in\Z_7\}.
      \]
      Jadi, $|X|=343$, $|H|=49$, dan $|V|=42$.
      Definisikan relasi ekuivalensi $R$ melalui partisi
      \[
            X/R=\{C_0,C_1,\ldots,C_6,C_*\},\qquad
            C_b=\{(a,b,0):a\in\Z_7\},\quad C_*=X\setminus H.
      \]
      Kelas $C_b$ memuat semua vektor dengan koordinat kedua $b$ dan
      koordinat ketiga $0$. Sebagai contoh,
      \[
            C_0=\{(0,0,0),(1,0,0),\ldots,(6,0,0)\},\qquad
            C_1=\{(0,1,0),(1,1,0),\ldots,(6,1,0)\}.
      \]
      Jadi, $C_0,\ldots,C_6$ membagi seluruh $H$, yaitu semua vektor
      yang koordinat ketiganya $0$. Vektor-vektor di $X$ yang belum
      tercakup adalah yang koordinat ketiganya bukan $0$; semuanya
      dikelompokkan ke dalam satu kelas tambahan
      \[
            C_*=\{(a,b,c):a,b\in\Z_7,\ c\in\{1,2,3,4,5,6\}\}.
      \]
      Misalnya, $(0,0,1)$, $(2,5,3)$, dan $(6,6,6)$ berada di $C_*$.
      Tanda $*$ hanya merupakan label kelas. Kelas ini memiliki
      $7\cdot7\cdot6=294$ anggota dan melengkapi partisi seluruh $X$.

      Setiap $C_b$ beririsan dengan $V$ pada enam anggotanya.
      Sebaliknya, $C_*\cap V=\varnothing$ karena koordinat ketiga
      setiap anggota $V$ adalah $0$. Oleh karena itu, $C_*$ tidak
      masuk aproksimasi atas $V$, sehingga
      \[
            \upp{V}=\bigcup_{b\in\Z_7}C_b=H,
            \qquad V\subsetneq\upp{V}\subsetneq X.
      \]
      Kita periksa bahwa $V$ merupakan ruang vektor rough atas $\Z_7$.
      \begin{enumerate}
            \item Untuk $u=(a,b,0)$ dan $v=(c,d,0)$ dalam $V$,
                  \[
                        u+v=(a+c,b+d,0)\in H=\upp{V}.
                  \]
                  Hasilnya tidak harus berada di $V$; misalnya,
                  $(1,2,0)+(6,3,0)=(0,5,0)\notin V$.
            \item Penjumlahan pada $\upp{V}$ bersifat asosiatif dan
                  komutatif karena diwarisi dari $X$.
            \item Vektor $0=(0,0,0)\in\upp{V}$ merupakan identitas.
                  Setiap $v=(a,b,0)\in V$ memiliki invers
                  $-v=(-a,-b,0)\in V$, sebab $a\ne0$ mengakibatkan
                  $-a\ne0$ di $\Z_7$.
            \item Untuk setiap $\lambda\in\Z_7$ dan $v=(a,b,0)\in V$,
                  \[
                        \lambda v=(\lambda a,\lambda b,0)\in H=\upp{V}.
                  \]
                  Kedua sifat distributif, kesesuaian perkalian skalar,
                  dan $1v=v$ berlaku karena operasinya diwarisi dari
                  ruang vektor $X$ atas $\Z_7$.
      \end{enumerate}
      Jadi, $V$ merupakan ruang vektor rough dengan aproksimasi atas
      yang bukan seluruh $X$. Namun, $V$ bukan ruang vektor biasa
      karena $0\notin V$. Selain itu, $\lowp{V}=\varnothing$ karena
      setiap $C_b$ memuat $(0,b,0)\notin V$.

      Jika $R$ diganti dengan relasi kesamaan $R_1$, maka
      $\overline{App}_{R_1}(V)=V$. Hasil penjumlahan
      $(1,2,0)+(6,3,0)=(0,5,0)\notin V$ menunjukkan bahwa
      $V$ bukan ruang vektor rough terhadap $(X,R_1)$.
\end{contoh}

\begin{teorema}\label{thm:21}
      Jika $V$ merupakan ruang vektor rough dan $V=\upp{V}$, maka $V$
      juga merupakan ruang vektor.
\end{teorema}
\begin{proof}
      Pembuktiannya langsung. Berdasarkan Definisi~\ref{def:15}, $V$
      memenuhi syarat-syarat ruang vektor.
\end{proof}

\begin{catatan}\label{rem:22}
      Irisan dua ruang vektor rough tidak harus merupakan subruang vektor
      rough. Sebagai contoh, misalkan $U=\Z_6\times\Z_6$ dan klasifikasi
      $U$ adalah $U/R=\{C_1,C_2\}$, dengan $C_1$ terdiri atas semua elemen
      yang sekurang-kurangnya satu komponennya nol dan $C_2=U\setminus C_1$.
      Misalkan
      \[
            V=\{(1,1),(0,2),(0,4),(5,5)\},\qquad
            W=\{(1,1),(5,5),(3,3),(4,0),(2,0)\}.
      \]
      Maka $V$ dan $W$ keduanya merupakan ruang vektor rough, sedangkan
      irisannya, $V\cap W=\{(1,1),(5,5)\}$, dengan aproksimasi
      $\upp{V\cap W}=C_2$, bukan merupakan ruang vektor rough karena
      $\upp{V\cap W}=C_2$ tidak memiliki identitas rough.
\end{catatan}

Pada proposisi berikut, diberikan syarat agar irisan dua ruang vektor
rough merupakan ruang vektor rough.

\begin{proposisi}\label{prop:23}
      Jika $(U,R)$ merupakan ruang aproksimasi dan $V$ serta $W$ merupakan
      dua ruang vektor rough atas lapangan $\F$ dengan operasi biner $(+)$
      dan perkalian skalar $(\cdot)$, maka $V\cap W$ juga merupakan subruang
      vektor rough atas lapangan yang sama apabila
      \begin{equation}
            \bigl[\upp{V}\cap\upp{W}\bigr]\subset\upp{V\cap W}.
            \tag{3.1}\label{eq:3.1}
      \end{equation}
\end{proposisi}
\begin{proof}
      Untuk menunjukkan bahwa $V\cap W$ merupakan subruang vektor rough,
      kita harus membuktikan bahwa $V\cap W$ merupakan subgrup rough dan
      $\lambda x\in\upp{V\cap W}$ untuk setiap $\lambda\in\F$ dan
      $x\in V\cap W$. Misalkan $x$ dan $y$ sebarang elemen dalam
      $V\cap W$, dan $\lambda\in\F$ sebarang. Maka $x+y$ berada di
      $\upp{V}$ dan $\upp{W}$ karena $V$ dan $W$ merupakan ruang vektor
      rough. Dengan demikian, $x+y\in\upp{V}\cap\upp{W}$, dan berdasarkan
      \eqref{eq:3.1}, diperoleh $x+y\in\upp{V\cap W}$. Selain itu,
      $x^{-1}\in V\cap W$, sehingga berdasarkan Definisi~\ref{def:10},
      $V\cap W$ merupakan subgrup rough.

      Sekarang, jika $x\in V\cap W$ dan $\lambda\in\F$, maka
      $\lambda x$ berada di $\upp{V}$ dan $\upp{W}$. Jadi,
      $\lambda x\in\upp{V}\cap\upp{W}$, dan kembali menggunakan asumsi
      \eqref{eq:3.1}, diperoleh $\lambda x\in\upp{V\cap W}$. Oleh karena
      itu, berdasarkan Definisi~\ref{def:17}, $V\cap W$ merupakan subruang
      vektor rough.
\end{proof}

\section{Struktur Aljabar}

\begin{definisi}\label{def:24}
      Misalkan $A=(U,R)$ suatu ruang aproksimasi. Ruang vektor rough $V$
      disebut \textbf{jumlah langsung} dari subruang-subruang vektor rough
      $W_1$ dan $W_2$ apabila
      \begin{enumerate}
            \item $\upp{W_1}\cap\upp{W_2}=\{0\}$, dan
            \item setiap elemen $x$ di $V$ dapat dituliskan sebagai $x=w_1+w_2$,
                  dengan $w_1\in\upp{W_1}$ dan $w_2\in\upp{W_2}$.
      \end{enumerate}
\end{definisi}
Perhatikan bahwa pada syarat pertama, $0$ merupakan identitas rough
dari grup rough $(V,+)$.

\begin{definisi}\label{def:25}
      Misalkan $A=(U,R)$ suatu ruang aproksimasi dan $V$ serta $W$ merupakan
      dua ruang vektor rough pada $A$ atas lapangan $\F$ yang sama.
      Pemetaan $L:\upp{V}\to\upp{W}$ disebut \textbf{linear rough}
      apabila berlaku:
      \begin{enumerate}
            \item $L(u+v)=L(u)*L(v)$; dan
            \item $L(\alpha v)=\alpha L(v)$,
      \end{enumerate}
      di mana $+$ merupakan operasi biner yang didefinisikan pada $V$ dan
      $*$ merupakan operasi biner yang didefinisikan pada $W$.
\end{definisi}

\begin{definisi}\label{def:26}
      Pemetaan linear rough disebut \textbf{isomorfisme rough} apabila
      pemetaan tersebut bijektif.
\end{definisi}

\begin{definisi}\label{def:27}
      Misalkan $A=(U,R)$ suatu ruang aproksimasi dan
      $L:\upp{V}\to\upp{W}$ suatu pemetaan linear rough. Maka himpunan
      \[
            RKer(L)=\{v\in\upp{V}:L(v)=_{\upp{W}}\}
            =L^{-1}\bigl(0_{\upp{W}}\bigr)
      \]
      disebut \textbf{kernel rough} dari $L$, dengan $_{\upr{W}}$
      merupakan identitas rough dari $W$.
\end{definisi}

\begin{catatan}\label{rem:28}
      Jika $L:\upp{V}\to\upp{W}$ merupakan isomorfisme rough, maka
      $L^{-1}:\upp{W}\to\upp{V}$ juga merupakan isomorfisme rough.
      Hal ini karena sifat linearitas berlaku untuk invers isomorfisme
      rough $L^{-1}$.
\end{catatan}

\begin{definisi}\label{def:29}
      Misalkan $V$ merupakan ruang vektor rough atas lapangan $\F$ dalam
      ruang aproksimasi $(U,R)$. Misalkan $N$ suatu subruang rough dari $V$.
      \textbf{Hasil bagi rough} dari $V$ oleh $N$ adalah himpunan
      \[
            V/N=\{x+\upp{N}:x\in\upp{V}\},
      \]
      di mana $x+N$ adalah himpunan semua $y$ yang memenuhi relasi
      $x-y\in\upp{N}$.
\end{definisi}

\begin{definisi}\label{def:30}
      Misalkan $V$ merupakan ruang vektor rough dalam ruang aproksimasi
      $(U,R)$ atas lapangan $\F$, dengan operasi biner $(+)$. Misalkan
      $N$ suatu subruang rough dari $V$. Maka himpunan
      \[
            x+\upp{N}=\{x+n:n\in\upp{N}\}
      \]
      disebut \textbf{koset rough} dari $N$ dalam $V$ yang memuat $x$,
      dengan $x$ suatu elemen tetap di $\upp{V}$.
\end{definisi}

Perhatikan bahwa, berdasarkan sifat komutatif pada $\upp{V}$,
diperoleh
\[
      x+\upp{N}=\upp{N};\qquad\forall x\in\upp{V},
\]
dan tidak ada perbedaan antara kedua bentuk tersebut.

\begin{teorema}\label{thm:31}
      Jika $V/N$ merupakan hasil bagi rough dari $V$ oleh $N$, maka
      \[
            \upp{V}=\bigcup V/N.
      \]
\end{teorema}
\begin{proof}
      Jelas bahwa $\bigcup V/N\subseteq\upp{V}$. Kita hanya perlu
      membuktikan bahwa $\upp{V}\subseteq\bigcup V/N$. Misalkan
      $x\in\upp{V}$ sebarang. Maka untuk setiap $N\subseteq V$,
      terdapat $0\in\upp{N}$ sehingga $x=x+0\in x+\upp{N}$. Jadi,
      $\upp{V}\subseteq\bigcup V/N$.
\end{proof}

\begin{lema}[Sifat-Sifat Koset Rough]\label{lem:32}
      Misalkan $N$ suatu subruang rough dari ruang vektor rough $V$ dan
      $x,y\in\upp{V}$.
      \begin{enumerate}
            \item $x\in x+\upp{N}$.
            \item $x+\upp{N}=\upp{N}$ jika dan hanya jika $x\in\upp{N}$.
            \item $x+\upp{N}=y+\upp{N}$ jika dan hanya jika $x\in y+\upp{N}$.
            \item $x+\upp{N}=y+\upp{N}$ atau
                  $\bigl(x+\upp{N}\bigr)\cap\bigl(y+\upp{N}\bigr)=\varnothing$.
            \item $x+\upp{N}=y+\upp{N}$ jika dan hanya jika $x-y\in\upp{N}$.
            \item $|x+\upp{N}|=|y+\upp{N}|$ apabila $V$ hingga.
            \item $x+\upp{N}$ merupakan subruang rough dari $V$ jika dan hanya
                  jika $x\in\upp{N}$.
      \end{enumerate}
\end{lema}
\begin{proof}
      \leavevmode
      \begin{enumerate}
            \item Karena $0\in\upp{N}$, maka $x=x+0\in x+\upp{N}$.

            \item Andaikan $x+\upp{N}=\upp{N}$. Berdasarkan sifat sebelumnya,
                  $x\in x+\upp{N}=\upp{N}$. Sebaliknya, andaikan $x\in\upp{N}$.
                  Karena $\upp{N}$ merupakan ruang vektor rough, maka
                  $x+\upp{N}=\upp{N}$.

            \item Andaikan $x+\upp{N}=y+\upp{N}$. Kita mengklaim bahwa
                  $x\in y+\upp{N}$. Berdasarkan sifat pertama,
                  $x\in x+\upp{N}=y+\upp{N}$. Sebaliknya, andaikan
                  $x\in y+\upp{N}$, dan kita mengklaim bahwa
                  $x+\upp{N}=y+\upp{N}$. Berdasarkan asumsi, $x=y+n$ dengan
                  $n\in\upp{N}$, sehingga
                  \[
                        x+\upp{N}=(y+n)+\upp{N}
                        =y+\bigl(n+\upp{N}\bigr)=y+\upp{N}.
                  \]
                  Dengan demikian, kesamaan tersebut berlaku.

            \item Andaikan terdapat elemen
                  $z\in\bigl(x+\upp{N}\bigr)\cap\bigl(y+\upp{N}\bigr)$.
                  Maka $z\in x+\upp{N}$ dan $z\in y+\upp{N}$. Berdasarkan sifat
                  ketiga, diperoleh $z+\upp{N}=x+\upp{N}=y+\upp{N}$.
                  Jadi, setiap dua koset sama atau saling lepas.

            \item $x+\upp{N}=y+\upp{N}$ jika dan hanya jika
                  $(x-y)+\upp{N}=0+\upp{N}$, jika dan hanya jika
                  $(x-y)\in0+\upp{N}=\upp{N}$.

            \item Untuk menunjukkan sifat ini, kita harus menemukan suatu
                  bijeksi dari $x+\upp{N}$ pada $y+\upp{N}$. Jelas bahwa fungsi
                  \[
                        x+n\longmapsto y+n;\qquad n\in\upp{N}
                  \]
                  memetakan $x+\upp{N}$ pada $y+\upp{N}$. Jika $y+n_1=y+n_2$,
                  maka $n_1=n_2$, sehingga fungsi tersebut satu-satu. Jadi,
                  $|x+\upp{N}|=|y+\upp{N}|$.

            \item Andaikan $x+\upp{N}$ merupakan subruang rough dari $V$.
                  Karena $x=x+0\in x+\upp{N}$, $0\in\upp{N}$, dan
                  $\upp{N}$ merupakan subruang rough, maka berdasarkan
                  Definisi~\ref{def:15}, diperoleh $x=x+0\in\upp{N}$.
                  Sebaliknya, jika $x\in\upp{N}$, maka berdasarkan sifat kedua,
                  diperoleh $x+\upp{N}=\upp{N}$. Himpunan $\upp{N}$ sudah
                  merupakan subruang rough, sehingga $x+\upp{N}$ juga demikian.
      \end{enumerate}
\end{proof}

\begin{definisi}\label{def:33}
      Misalkan $V$ merupakan ruang vektor rough dalam ruang aproksimasi
      $(U,R)$. Himpunan bagian $F\subset\upp{V}$ disebut
      \textbf{simetris} apabila untuk setiap $x\in F$, berlaku $(-x)\in F$.
\end{definisi}

\begin{akibat}\label{cor:34}
      Jika $V$ merupakan ruang vektor rough dalam ruang aproksimasi $(U,R)$,
      maka $V$ dan setiap subruang rough dari $V$ merupakan himpunan simetris.
\end{akibat}

\begin{thebibliography}{9}
      \bibitem{bagirmaz}
      N.~Ba\u{g}{\i}rmaz, \.{I}\c{c}en, \.{I}., A.~F.~\"Ozcan,
      \textit{Topological Rough Groups},
      Topological Algebra and its Applications, 4(1), 2016, 31--38.
      \url{https://doi.org/10.1515/taa-2016-0004}.

      \bibitem{biwa}
      R.~Biwa, S.~Nanda,
      \textit{Rough Groups and Rough Subgroups},
      Bulletin of the Polish Academy of Sciences--Mathematics,
      42(3), 1994, 251--254.

      \bibitem{davvaz}
      B.~Davvaz, \textit{Roughness in rings},
      Information Sciences, 164(1--4), 2004, 147--163.
      \url{https://doi.org/10.1016/j.ins.2003.10.001}.

      \bibitem{modules}
      B.~Davvaz, M.~Mahdavipour, \textit{Roughness in modules},
      Information Sciences, 176(24), 2006, 3658--3674.
      \url{https://doi.org/10.1016/j.ins.2006.02.014}.

      \bibitem{pawlak}
      Z.~Pawlak, \textit{Rough sets},
      International Journal of Computer \& Information Sciences,
      11(5), 1982, 341--356.
\end{thebibliography}

\end{document}
